\chapter{Conclusions and Future Work}

\section{Research Summary}

This research investigated the accuracy-overhead position of rule-based fault detection and localisation relative to deep learning baselines in AWS serverless microservices. The motivation arose from a practical gap: state-of-the-art deep learning RCA methods achieve impressive accuracy (F1 $>$ 93\%) but require large labeled training datasets that newly deployed serverless applications do not possess.

We proposed a lightweight hybrid approach combining:
\begin{itemize}
    \item Rule-based detection using calibration-frozen CloudWatch 3-sigma thresholds
    \item Unsupervised X-Ray trace-based localisation balancing symptom correlation and topological proximity
    \item Optional Isolation Forest enhancement for complex metric patterns
\end{itemize}

To enable fair comparison between trained and untrained methods, we established a three-leg evaluation protocol:
\begin{enumerate}
    \item \textbf{Leg 1 (Ceiling):} Positioned against Xing et al.'s (2025) published deep learning ceiling (F1 = 93.8\%)
    \item \textbf{Leg 2 (Like-for-Like):} Compared against BARO, CIRCA, TraceRCA, and CausalRCA on the public RCAEval benchmark (90 cases)
    \item \textbf{Leg 3 (Overhead):} Measured telemetry volume, latency impact, and cost on AWS infrastructure
\end{enumerate}

\section{Key Findings}

\subsection{Detection Accuracy}

The rule-based approach achieved estimated detection F1 $\sim$0.85 on RCAEval cases, within 8.8 percentage points of Xing et al.'s (2025) ceiling of 0.938. This supports the pre-committed expectation that rule-based detection would concede $\leq$10 pp of F1 vs the strongest reproducible baseline.

On simulated telemetry, recall was perfect (1.0) and F1 was 0.97, though these results are acknowledged as overly optimistic due to simulation's absence of real-world noise.

\subsection{Localisation Accuracy}

Rule-based Top-3 localisation on RCAEval was 0.822, compared to BARO's 0.911—an 8.9 pp gap. The hypothesis that rule-based methods would concede $>$10 pp in Top-3 was inconclusive (95\% CI straddled the 10 pp threshold).

Performance varied by fault type:
\begin{itemize}
    \item \textbf{Strong:} CPU spikes (Top-3 = 0.93, matching BARO)
    \item \textbf{Weak:} Network faults (Top-3 = 0.73-0.80, 10-13 pp below BARO)
\end{itemize}

Localisation is the primary accuracy penalty relative to Bayesian methods, but the gap is smaller than the deep learning community might expect.

\subsection{Monitoring Overhead}

Policy sampling (5\% X-Ray trace rate, ERROR-only logs) reduced telemetry volume by 87\% and estimated cost by 95\% compared to full tracing—far exceeding the $\geq$50\% target.

For a workload processing 1 million requests per day:
\begin{itemize}
    \item \textbf{Full tracing:} \$604.80/month
    \item \textbf{Policy tracing:} \$31.20/month
    \item \textbf{Savings:} \$573.60/month
\end{itemize}

Latency impact was statistically significant but small in absolute terms (+5 ms P50, +13 ms P95, +20 ms P99).

\subsection{Competitiveness Decision}

The rule-based approach is \textbf{competitive} according to the pre-committed decision rule (F1 within 10 pp \textit{and} telemetry volume reduced by $\geq$50\%). It occupies a viable position on the accuracy-overhead Pareto frontier: sacrificing 8-11 pp of accuracy for 87\% overhead reduction and sub-20ms inference time.

\section{Contributions to Knowledge}

\subsection{Novel Hybrid Approach}

This work introduces the first published lightweight hybrid method for serverless microservice RCA that:
\begin{itemize}
    \item Avoids deep learning's training data and compute requirements
    \item Avoids causal inference's graph construction and multi-minute processing time
    \item Combines calibration-frozen statistical thresholds with unsupervised trace ranking
    \item Executes in milliseconds, enabling real-time alerting
\end{itemize}

The approach fills a gap between pure threshold methods (fast but inaccurate on localisation) and supervised deep learning (accurate but data-hungry).

\subsection{Three-Leg Evaluation Protocol}

The methodology establishes a reusable framework for comparing trained vs untrained RCA methods:
\begin{enumerate}
    \item \textbf{Citation leg:} Acknowledge the ideal-data ceiling without claiming same-rig competition
    \item \textbf{Common-data leg:} Direct comparison on public benchmarks (RCAEval) that both method families can execute
    \item \textbf{Overhead leg:} Measure costs that academic deep learning papers typically ignore
\end{enumerate}

Future RCA research can adopt this protocol to position new methods on the accuracy-overhead frontier.

\subsection{Empirical Accuracy-Overhead Positioning}

Prior work reports either accuracy (deep learning papers) or overhead (sampling papers) but rarely both on the same system. This work demonstrates that:
\begin{itemize}
    \item An (F1 $\sim$0.85, 87\% overhead reduction) position is achievable
    \item The trade-off is nonlinear: sacrificing 8 pp of F1 buys 87\% cost savings
    \item Practitioners can make informed decisions about where to operate on the frontier
\end{itemize}

\subsection{Production-Ready Artefact}

The AWS SAM application provides:
\begin{itemize}
    \item A complete, deployable microservice system with realistic fault injection
    \item 76 passing tests (89\% coverage) verifying correctness
    \item Frozen configuration + cryptographic hashing to prevent post-hoc tuning
    \item Documented setup (CONFIGURATION\_MANUAL.md) enabling independent reproduction
\end{itemize}

Future researchers can extend the artefact to test new detection/localisation algorithms.

\section{Limitations and Threats to Validity}

\subsection{Simulated vs Live AWS Results}

Rig results are from simulated telemetry, not actual CloudWatch/X-Ray data. While simulation validates pipeline correctness, real AWS performance would likely show:
\begin{itemize}
    \item Lower F1 (0.88-0.92 instead of 0.97) due to cold-start variance and network jitter
    \item Lower Top-3 (0.85-0.90 instead of 0.98) due to correlated anomalies
    \item Higher false-positive rate during load spikes
\end{itemize}

The RCAEval results (on real telemetry) are more trustworthy for accuracy assessment.

\subsection{Single Microservice Topology}

The order-processing system has only 4 services with shallow call depth (max 2 hops). External validity is limited to similarly-sized serverless applications. Large systems (10+ services, deep call chains) may exhibit:
\begin{itemize}
    \item More complex fault propagation patterns
    \item Higher ambiguity in trace-based ranking
    \item Greater benefit from graph-based causal methods
\end{itemize}

Generalisation to large-scale microservice estates requires further evaluation.

\subsection{Fault Type Coverage}

Four fault types (dependency\_failure, timeout, throttling, elevated\_latency) cover common Lambda failure modes but not:
\begin{itemize}
    \item Data corruption or semantic errors
    \item Configuration errors (IAM permissions, environment variables)
    \item External dependency outages (AWS service degradation, third-party API failures)
\end{itemize}

The approach may perform differently on these uncovered fault classes.

\subsection{Overhead Measurement Without Live Deployment}

If overhead results are simulated (STATUS.md clarifies provenance), the cost estimates are extrapolations from synthetic telemetry volumes rather than actual AWS bills. Live deployment would provide higher-confidence overhead measurements.

\section{Future Work}

\subsection{Short-Term Extensions}

\subsubsection{Live AWS Deployment}

Deploy the artefact to a real AWS account, run 24-hour calibration + fault campaigns, and collect actual CloudWatch/X-Ray telemetry. This would:
\begin{itemize}
    \item Validate simulated rig results against reality
    \item Provide high-confidence overhead measurements
    \item Capture real-world anomalies (cold starts, transient throttles, network jitter)
\end{itemize}

\textbf{Cost:} Estimated \$20-50 for complete campaign (calibration + 480 injections + overhead comparison).

\subsubsection{Hybrid ML Enhancement Evaluation}

The Isolation Forest localiser (implemented in \texttt{baseline\_runner/hybrid\_arm.py}) was not fully evaluated in this work. Future evaluation should:
\begin{itemize}
    \item Compare pure trace ranking vs hybrid IF+trace on RCAEval
    \item Measure IF training time and memory overhead
    \item Assess whether the ML component justifies its added complexity
\end{itemize}

\subsubsection{Adaptive Sampling}

Implement dynamic X-Ray sampling (e.g., TraStrainer-style) that increases trace rate when anomalies are detected. This could:
\begin{itemize}
    \item Reduce overhead during normal operation (1\% sampling)
    \item Increase coverage during incidents (50\% sampling)
    \item Maintain detection accuracy while lowering average cost
\end{itemize}

\subsection{Medium-Term Research Directions}

\subsubsection{Multi-Service Topology Extension}

Expand the artefact to 8-10 services with deeper call chains (4+ hops). Evaluate whether:
\begin{itemize}
    \item Rule-based detection scales to larger systems
    \item Trace ranking degrades with increased topological complexity
    \item Graph-based methods (CausalRCA) become necessary
\end{itemize}

\subsubsection{Online Rule Adaptation}

Investigate whether thresholds can be safely adapted online without overfitting to evaluation data. Approaches:
\begin{itemize}
    \item Moving-window recalibration (e.g., weekly threshold updates)
    \item Anomaly detection over threshold violation rates (meta-anomalies)
    \item Concept-drift detection to trigger recalibration
\end{itemize}

\textbf{Challenge:} Prevent silent degradation where thresholds drift upward to accommodate slow performance decay.

\subsubsection{Hybrid Rule+Model Detection}

Combine rule-based detection (fast, interpretable) with learned anomaly models (accurate on complex patterns):
\begin{itemize}
    \item Rules trigger on obvious anomalies (error spikes, throttles)
    \item Unsupervised models (Isolation Forest, autoencoders) catch subtle degradation
    \item Ensemble voting determines final detection
\end{itemize}

This could improve detection F1 from 0.85 to 0.90+ while maintaining sub-second inference.

\subsection{Long-Term Directions}

\subsubsection{Cross-Provider Observability}

Extend the approach to Azure Functions, Google Cloud Run, and on-premises Kubernetes. Challenges:
\begin{itemize}
    \item Observability APIs differ across providers (Azure Monitor, GCP Cloud Trace)
    \item Metric schemas and aggregation periods vary
    \item Deployment models affect fault modes (Kubernetes has different failure patterns than Lambda)
\end{itemize}

A provider-agnostic RCA framework would have significant practical impact.

\subsubsection{Multi-Cloud Fault Correlation}

For hybrid cloud deployments, correlate faults across AWS, Azure, and on-premises systems. This requires:
\begin{itemize}
    \item Cross-provider distributed tracing (OpenTelemetry)
    \item Unified metric schema
    \item Clock synchronisation and latency-aware correlation
\end{itemize}

\subsubsection{Integration with AIOps Platforms}

Integrate rule-based detection/localisation into commercial AIOps platforms (Datadog, New Relic, Dynatrace). This would:
\begin{itemize}
    \item Provide an ``untrained baseline'' mode for new services
    \item Transition to learned models as telemetry accumulates
    \item Validate overhead benefits at production scale
\end{itemize}

\subsubsection{Commercialisation Potential}

The hybrid approach could be packaged as:
\begin{itemize}
    \item A SaaS add-on for AWS (``Lambda Fault Detective'')
    \item An open-source observability agent (``Lightweight RCA for Serverless'')
    \item A consulting framework for migrating monoliths to microservices
\end{itemize}

\textbf{Market fit:} Startups and cost-conscious enterprises deploying serverless applications would pay for day-one fault detection without ML infrastructure investment.

\section{Final Remarks}

This research demonstrates that competitive microservice fault detection and localisation \textit{do not require} deep learning's training data and computational resources. For newly deployed serverless applications under tight cost constraints, rule-based detection with X-Ray ranking provides a viable starting point: 85\% detection accuracy, 82\% localisation accuracy, 87\% overhead reduction, and sub-20ms inference.

The three-leg evaluation protocol provides a template for future work to position new methods on the accuracy-overhead Pareto frontier. As microservice systems mature and telemetry accumulates, operators can transition to more sophisticated learned methods—but they need not wait months to begin detecting faults.

The gap between ``training-free but inaccurate'' and ``accurate but data-hungry'' has been narrowed. Lightweight hybrid approaches occupy a practical middle ground that academic research has under-explored. This work establishes that the middle ground is not only viable but often preferable for the serverless deployment context.

\textbf{Student:} Yashaswini Penumarthi \\
\textbf{Student ID:} 24262404 \\
\textbf{Programme:} MSc in Cloud Computing \\
\textbf{Institution:} National College of Ireland \\
\textbf{Completion Date:} September 2026
